\documentclass[12pt]{article}
\usepackage{hyperref}

\usepackage{graphicx}

\usepackage{mathtools}
\DeclarePairedDelimiter\ceil{\lceil}{\rceil}
\DeclarePairedDelimiter\floor{\lfloor}{\rfloor}




\usepackage{color}
\definecolor{heavyblue}{cmyk}{1,1,0,0.25}

\usepackage{amsmath}
\usepackage{amsfonts}
\usepackage{amssymb}

\usepackage{listings}
\lstset{language=C++,
                basicstyle=\ttfamily,
                keywordstyle=\color{blue}\ttfamily,
                stringstyle=\color{red}\ttfamily,
                commentstyle=\color{green}\ttfamily,
                morecomment=[l][\color{magenta}]{\#}
}

\usepackage{breqn}
\usepackage{physics}

\newcommand\fnurl[2]{%
  \href{#2}{#1}\footnote{\url{#2}}%
}

\newcommand{\Four}{\mathcal{F}}

% these bracket shortcuts break algorithm packages:
%\renewcommand{\(}{\left(}
%\renewcommand{\)}{\right)}
%% \renewcommand{\{}{\left{}
%% \renewcommand{\}}{\right}}

\newcommand{\thetitle}{Real-complex multi-dimensional FFTs with transposition to use the even-length method FIXME: this is a bad title.}
\newcommand{\theauthor}{Malcolm Roberts \\
AIG Software Development \\
malcolm.roberts@amd.com}


\hypersetup{
  pdftitle={\thetitle},
  pdfpagemode=UseOutlines,
  citebordercolor=0 0 1,
  colorlinks=true,
  allcolors=heavyblue,
  breaklinks=true,
  pdfauthor={\theauthor},
  pdfpagetransition=Dissolve,
  bookmarks=true
}

\title{\thetitle}
\author{\theauthor}


\def\no#1{\nobreak{#1}} % stop equation breaking in dmath

\usepackage{algorithm}
\usepackage{algpseudocode}

\begin{document}
\maketitle

\section{Introduction}

Fast Fourier transforms (FFTs) are an important algorithm in
high-performance computing, offering improved numerical accuracy and
reduced computational and bandwidth costs.  While the Fourier
transform is most naturally defined on complex values, applications of
the Fourier transform generally use real values.  This document
outlines a new method for efficiently computing real-to-complex FFTs in
situations where conventional methods are not applicable.  This
method, which we refer to as the real-transpose method, allows for
increased performance and accuracy compared to alternatives.

FIXME: add N log N and numerical error.

\section{Problem statement}

The FFT of the one-dimensional input $x=\{x_n\}_{n=0}^{N-1}$ is 
\begin{dmath}
  \label{dft}
  X_k = \mathcal{F}\left[x\right]_k = \sum_{n=0}^{N-1} \omega_N^{kn}x_n
\end{dmath}
for $k=0,\dots,{N-1}$, and where $\omega_N=e^{\frac{-2\pi{}i}{N}}$ is
the $N$th root of unity.  If $x$ is real-valued, then
equation~\eqref{dft} exhibits Hermitian symmetry.  That is,
\begin{dmath}
  X_k = X^*_{-k}
\end{dmath}
where ${}^*$ denotes complex conjugation.  As such, we can save flops
and bytes by not computing redundant values.  In particular, we
store the following $\floor{N/2}+1$ values
\begin{dmath}
  X_0, X_1, \dots, X_{\floor{N/2}}
\end{dmath}
and note that $X_0$ is guaranteed to be real-valued, as is
$X_{\floor{N/2}}$ if $N$ is even.

\subsection{Complex embedding}
\label{sscomplexembed}

The simplest way to compute a real-to-complex FFT is to embed the real
data into a complex buffer with zero real part, perform a
complex-to-complex FFT of length $N$, and then copy the first
$\floor{N/2}+1$ values into the output buffer.  This is outline in
alrogithm~\ref{complexembed1d}.  The memory footprint is roughly twice
the size of the real-valued input, and the underlying
complex-to-complex FFT is of length $N$..  The inverse,
complex-to-real transform follows much as one would expect.  This
algorithm works for general problem sizes.
\begin{algorithm}
  \caption{Complex embedding}\label{complexembed1d}
  \begin{algorithmic}[1]
    \For{$n\gets 0, N-1$}
    \State $z_n \gets x_n$
    \EndFor
    \State $Z \gets \Four_N\left[z\right]$
    \For{$n\gets 0, \floor{N/2}$}
    \State $X_n \gets Z_n$
    \EndFor
    \State \textbf{return} $X$
  \end{algorithmic}
\end{algorithm}

\subsection{Paired transforms}

Improvements to the complex embedding method rely, more-or-less, on
the fact that there are two real values in each complex value.  This
is particularly clear in the paired-transform method.  If we have two
real-valued inputs $x$ and $y$ which we wish to transform, we can make
use of the linearity of the Fourier transform to compute the transform
with a lower cost than that in subsection~\ref{sscomplexembed}.  In
particular, let $z_n = x_n + i y_n$.  Then, $Z_k = X_k + i Y_k$, and
we can recover the invididual transforms via
\begin{align}
  X_k &= \frac{Z_k + Z_{N-k}^*}{2} \\
  Y_k &= \frac{Z_k - Z_{N-k}^*}{2i}.
\end{align}
This entails one length-$N$ complex-to-complex FFT for two inputs,
which has the same memory footprint as the two real-valued inputs.
The algorithm is described in algorithm~\ref{paired}.  This method
works for both even and off $N$, but requires an even number of
transforms.  If $x$ and $y$ have very different magnitudes, then this
method is not numerically robust, as the smaller input will have a
numerical error dominated by the larger input.
\begin{algorithm}
  \caption{Paired transforms}\label{paired}
  \begin{algorithmic}[1]
    \For{$n\gets 0, N-1$}
    \State $z_n \gets x_n + i y_n$
    \EndFor
    \State $Z \gets \Four_N\left[z\right]$
    \For{$k\gets 0, \floor{N/2}$}
    \State $X_k \gets \frac{Z_k + Z_{N-k}}{2}$
    \State $Y_k \gets \frac{Z_k - Z_{N-k}}{2i}$
    \EndFor
    \State \textbf{return} $X$, $Y$
  \end{algorithmic}
\end{algorithm}

\subsection{Even-length transforms}

When $N$ is even, we have the opportunity to use a more efficient
algorithm, where we treat the data as complex-valued from the
beginning.  In particular, let $z_n = x_{2n} + i x_{2n+1}$ for
$n=0,\dots,N/2-1$.  Then,
\begin{align}
  X_0 &= \Re{Z_0} - Im{Z_0}\\
  X_k &= \frac{1}{2} \left(Z_k + Z_{\frac{N}{2}-k}^*\right)
  + \omega_N^k \frac{i}{2} \left(Z_{\frac{N}{2}-k}^* - Z_{}k\right) \\
  X_{\frac{N}{2}} &=  \Re{Z_0} + Im{Z_0}
\end{align}
This entail a complex FFT of length $N/2$, and the memory footprint
for this transform is basically the same as that of the single input.
Since this method only relies on one input, it doesn't have the same
numerical robustness problems as the paired algorithm.
\begin{algorithm}
  \caption{Even algorithm}\label{even}
  \begin{algorithmic}[1]
    \For{$n\gets 0, N/2-1$}
    \State $z_n \gets x_{2n} + i x_{2n+1}$
    \EndFor
    \State $Z \gets \Four_{N/2}\left[z\right]$
    \State $X_0 \gets \Re{Z_0} - \Im{Z_0}$ 
    \For{$k\gets 1, N/2-1$}
    \State $X_k \gets \frac{1}{2} \left(Z_k + Z_{\frac{N}{2}-k}^*\right)
  + \omega_N^k \frac{i}{2} \left(Z_{\frac{N}{2}-k}^* - Z_{}k\right)$
    \EndFor
    \State $X_{\frac{N}{2}} \gets \Re{Z_0} + \Im{Z_0}$ 
    \State \textbf{return} $X$
  \end{algorithmic}
\end{algorithm}

The paired algorithm offers performance advantages, and is the
preferred algorithm when it is available.

\section{Multi-dimensional transforms}

Many applications make use of 2D or 3D FFTs.  The standard method for
compute a $d$-dimensional FFT with length
$N_1\times{}\dots{}\times{}N_d$ is to perform batched 1D transforms in
each dimension.  For real-to-complex transforms, the output exhibits
Hermitian symmetry, so only about half of the values are needed in the
output.

In the 1D case, there isn't much choice on how to apply this symmetry;
one keeps either the lower half or upper half of the array, and it's
just much simpler to keep the lower half.  In higher dimensions, the
symmetry is over some plane that bisects the data, and it saves a lot
of trouble to make this plane perpendicular to one of the axes.  In a
row-major format, uers expect $N_d$ to be chopped in half, and in a
column-major format, one expects $N_1$ to be chopped in half.  For
simplicity, let us stick to the row-major case for the remainder of
the discussion, keeping in mind that the results are applicable in
both cases.  Let $N_d' = \floor{N_d/2}+1$.


\begin{algorithm}
  \caption{Standard $D$-dimensional real-to-complex FFT}\label{DFFTold}
  \begin{algorithmic}[1]
    \For{$\left(n_1,\dots,N_{D-1}\right) \in \left(0,\dots,N_1-1\right) \times \left(0,\dots,N_{D-1}-1\right) $}
    \For{$n_D \in \left(0,\dots,N_D-1\right)$}
    \State $z_{n_1,n_2,\dots,n_D} \gets x_{n_1,n_2,\dots,n_D}$ \Comment{$z$ is complex}
    \EndFor
    \State $\left\{z_{n_1,n_2,\dots,n_d}\right\}_{n_1=0}^{N_1-1}
    \gets \Four_{N_1}\left[\left\{z_{n_1,n_2,\dots,n_d}\right\}_{n_1=0}^{N_1-1}\right]$
    \For{$n_1 \in \left(0,\dots,N_D'-1\right)$}
    \State $X_{n_1,n_2,\dots,n_d} \gets z_{n_1,n_2,\dots,n_d}$
    \EndFor
    \EndFor
    \For{$d \in \left(1, \dots, D-1\right)$}
    \For{$\left(n_1, n_2, \dots, n_{d-1}, n_{d+1}, \dots,n_D \right) \in
      \left(0,\dots, N_D'-1\right) \times
      \prod_{b\neq{}d} \left(0,\dots,N_b-1 \right)$}
    \State $\left\{X_{n_1,n_2,\dots,n_d}\right\}_{n_d=0}^{N_d-1}
    \gets \Four_{N_d}\left[\left\{X_{n_1,n_2,\dots,n_d}\right\}_{n_d=0}^{N_d-1}\right]$
    \EndFor
    \EndFor
    \State \textbf{return} $X$
  \end{algorithmic}
\end{algorithm}

When computing multi-dimensional real-to-complex FFTs, the first
batched 1D transforms take real data as input and produce complex
(Hermitian symmetric) data as output.  If $N_d$ is even, then we can
apply the even-length real/complex algorithm, which allows for
efficient, numerically robust computation.  If $N_d$ is not even, then
we woud fall back to less optimal algorithm.  However, if one of the
other dimensions is even, one is sorely tempted to use the even-length
method in that direction.  The result is indeed the Fourier transform
of the input, but with the symmetry taken in a different direction.
This isn't a problem for applications of the convolution theorem,
where one just performs a pointwise multiplication and then inverts
the transform.  But if the user needs to do a pointwise operation
which depends on the multi-index of the output, this will produce
incorrect results, since the output is transposed and half of it is
the complex conjugate of the expected result.

We can resolve this issue by re-applying the Hermitian symmetry
condition to recover the expected output.  Suppose $N_d$ is odd, and
$N_e$ is even.  Then we can compute the real-to-complex batched 1D
FFTs in the $N_e$ direction, producing a purely complex output,
perform the remaining batched 1D FFTs in the other directions.  The
Hermitian symmetry is the re-imposed by transposing the data and
taking the complex conjugate as approporiate.

The even-length method is about twice as efficient as the complex
embedding method, since the length of the underlying complex FFT is
half in the even case as in the embedded case.  If all of the
dimensions are roughly equal in magnitude, then the transpose-to-even
method has computational cost
\begin{dmath}
  \frac{k}{2} N_e \log{N_e} \prod_{d\neq e} N_d
  + \sum_{d \neq e}  k N_d \log{N_d}  \prod_{b\neq d} N_b
  =  \left(\frac{k}{2}  \log{N_e} + \sum_{d \neq e} k \log{N_d} \right)
  \prod_{d}{N_d}
\end{dmath}
For 2D problems, this is about 25\% less than the complex-embedding
method, and 17\% less for 3D problems.  In general, a (roughly
equiproportional) $D$-dimensional FFT would be sped up by a factor of
\begin{dmath}
  \frac{D - \frac{1}{2}}{D}.
\end{dmath}
Since smaller FFTs have better cache behaviour, the performance
improvement is likely better in practice.  For example, GPU-based 1D
FFTs using LDS will have more global-memory round-trips; performing a
length-$N$ vs a length $N/2$ complex FFT could conceivably not just
double but quadruple the bandwidth costs if the shorter transform fits
into LDS and the longer did not.

\begin{algorithm}
  \caption{$D$-dimensional real-to-complex FFT using the even-transpose method}\label{DFFTold}
  \begin{algorithmic}[1]
    \For{$\left(n_1,\dots,n_{e-1},n_{e+1}, \dots, n_{D}\right) \in \prod_{d\neq{}e}\left(0,\dots,N_d-1\right)$}
    \For{$n_e \in \left(0,\dots,N_e/2-1\right)$}
    \State $z_{n_e} \gets x_{n_1,\dots,2n_{e-1},\dots,n_D} + i x_{n_1,\dots,2n_{e-1}+1,\dots,n_D}$ 
    \EndFor
    \State $\left\{z_{n_e}\right\}_{n_e=0}^{N_e/2-1}
    \gets \Four_{N_e/2}\left[\left\{z_{n_e}\right\}_{n_1=0}^{N_e/2-1}\right]$
    \State $X_{n_1,\dots,n_{e-1},0,n_{e+1},\dots,n_D} \gets \Re{z_{0}} - \Im{z_{N_e-1}}$
    \For{$n_e \in \left(0,\dots,N_e'-1\right)$}
    \State $X_{n_1,\dots,n_D} \gets \frac{1}{2} \left(z_{n_e} + z_{\frac{N_e}{2}-k}^*\right)
    + \omega_{N_e}^k \frac{i}{2} \left(z_{\frac{N_e}{2}-k}^* - z_{n_e}\right)$
    \EndFor
    \State $X_{n_1,\dots,n_{e-1},N_e/2,n_{e+1},\dots,n_D}  \gets \Re{z_{0}} + \Im{z_{N_e-1}}$
    \EndFor
    \For{$d \in \left(1, \dots, e-1, e+1, \dots, D\right)$}
    \State $S =
      \prod_{\substack{b=0\\ b\notin{}\{d,e\}}}^{e-1} \left(0,\dots,N_b-1 \right)
      \times \left(0,\dots,N_e' \right)
      \times \prod_{\substack{b=e+1\\ b\notin{}\{d,e\}}}^{D} \left(0,\dots,N_b-1 \right)$
      \For{$\left(n_1, n_2, \dots, n_{d-1}, n_{d+1}, \dots,n_D \right) \in S$}
    \State $\left\{X_{n_1,n_2,\dots,n_d}\right\}_{n_d=0}^{N_d-1}
    \gets \Four_{N_d}\left[\left\{X_{n_1,n_2,\dots,n_d}\right\}_{n_d=0}^{N_d-1}\right]$
    \EndFor
    \EndFor
    \State $S =
      \prod_{d=1}^{e-1} \left(0,\dots,N_d-1 \right)
      \times \left(0,\dots,N_e' \right)
      \times \prod_{d=e+1}^{D} \left(0,\dots,N_d-1 \right)$
    \For{$\left(n_1, \dots, n_{e-1}, n_{e+1}, \dots,n_{D} \right) \in S$}
    \If{$n_e > N_e'$}
    \State $X_{n_1,\dots, N_e - n_e ,\dots, N_d - n_d} \gets X_{n_1,\dots, n_e,  \dots, n_D}^*$ \Comment{$X$ changes shape here}
    \EndIf
    \EndFor
    \State \textbf{return} $X$
  \end{algorithmic}
\end{algorithm}


FIXME: mention that this is novel

FIXME: algorithm.

\end{document}

